% !TeX program = lualatex
% !TeX encoding = UTF-8
% !TeX spellcheck = pt_BR

\documentclass[11pt,a4paper]{article}

\usepackage{fontspec}
\defaultfontfeatures{Ligatures=TeX}
% Mesma tipografia do relatorio.tex (Latin Modern Roman para o corpo do texto,
% Latin Modern Sans e Latin Modern Mono para títulos e código).
\setmainfont{Latin Modern Roman}
\setsansfont{Latin Modern Sans}
\setmonofont{Latin Modern Mono}

\usepackage[brazil]{babel}
\usepackage{geometry}

% Formato vertical estilo tela de smartphone (proporção 9:16, como a maioria
% dos celulares em retrato), para que cada página ocupe a tela inteira ao
% ser aberta e ampliada em um leitor de PDF no celular.
\geometry{
  paperwidth=90mm,
  paperheight=160mm,
  top=9mm,
  bottom=11mm,
  left=6mm,
  right=6mm
}

\usepackage{microtype}
\usepackage{xcolor}
\usepackage{graphicx}
\usepackage{booktabs}
\usepackage{tabularx}
\usepackage{enumitem}
\usepackage{amssymb}
\usepackage{tikz}
\usepackage{fancyhdr}
\usepackage{lastpage}
\usepackage{xurl}
\usepackage{hyperref}

% Mesma paleta de cores do relatorio.tex (linkblue, commentgreen, stringred,
% codegray, rulegray), para que os dois documentos pareçam parte do mesmo
% conjunto em vez de usarem identidades visuais distintas. "primaryColor" e
% "accentColor" são apelidos semânticos para o mesmo azul do relatório, e
% "greenAccent"/"redAccent" reaproveitam as cores já usadas lá para
% comentários/strings de código, aqui usadas para o LED verde/vermelho.
\definecolor{linkblue}{HTML}{1D4ED8}
\definecolor{codegray}{HTML}{F3F4F6}
\definecolor{rulegray}{HTML}{D1D5DB}
\definecolor{commentgreen}{HTML}{166534}
\definecolor{stringred}{HTML}{B91C1C}

\colorlet{primaryColor}{linkblue!85!black}
\colorlet{accentColor}{linkblue}
\colorlet{greenAccent}{commentgreen}
\colorlet{redAccent}{stringred}
\colorlet{cardBg}{codegray}
\colorlet{cardBorder}{rulegray}
\colorlet{mutedText}{black!45}

\hypersetup{
  colorlinks=true,
  linkcolor=accentColor,
  urlcolor=accentColor
}

% Estilo de cabeçalho e rodapé
\pagestyle{fancy}
\fancyhf{}
\renewcommand{\headrulewidth}{0pt}
\renewcommand{\footrulewidth}{0.4pt}

\lfoot{\tiny \color{mutedText} CESS-UFF / SGIMP}
\rfoot{\tiny \color{mutedText} Slide \thepage\ de \pageref{LastPage}}

\setlist[itemize]{leftmargin=1.2em,itemsep=0.3em,topsep=0.3em}
\setlist[enumerate]{leftmargin=1.4em,itemsep=0.3em,topsep=0.3em}
\setlength{\parindent}{0pt}

\newcommand{\WokwiURL}{https://wokwi.com/projects/471528241540407297}
\newcommand{\GitHubURL}{https://github.com/Argus-Kepheus/cess-uff}
\newcommand{\CandidateWebsite}{https://argus-kepheus.github.io/sergio-cordero-calvimontes/\#home}
\newcommand{\CandidateName}{\href{\CandidateWebsite}{Sergio Cordero Calvimontes}}
\newcommand{\okmark}{\textcolor{greenAccent}{\checkmark}}

% Título de slide
\newcommand{\slideheader}[1]{%
  {\Large\bfseries\color{primaryColor} #1}\par
  \vspace*{2mm}
  {\color{accentColor}\hrule height 1.2pt}
  \vspace*{4mm}
}

% Largura de card calculada a partir da largura útil da página corrente
% (\linewidth), em vez de um valor fixo em mm -- assim o card sempre ocupa a
% tela inteira disponível, mesmo que a geometria da página mude.
\newcommand{\cardtextwidth}{\dimexpr\linewidth-12pt\relax}

% Card padrão
\newcommand{\card}[2]{%
  \begin{tikzpicture}
    \node[fill=cardBg, draw=cardBorder, rounded corners=6pt, inner sep=6pt, text width=\cardtextwidth] {
      \small\textbf{\color{primaryColor} #1}\\[3pt]
      \footnotesize\color{black} #2
    };
  \end{tikzpicture}\par\vspace{3mm}
}

% Card de destaque, com cor de ênfase configurável
\newcommand{\alertcard}[3]{%
  \begin{tikzpicture}
    \node[fill=#1!8, draw=#1!40, rounded corners=6pt, inner sep=6pt, text width=\cardtextwidth] {
      \small\textbf{\color{#1} #2}\\[3pt]
      \footnotesize\color{black} #3
    };
  \end{tikzpicture}\par\vspace{3mm}
}

\begin{document}

% ------------------------------------------------------------------------------
% SLIDE 1: Capa
% ------------------------------------------------------------------------------
\begin{center}
  \vspace*{4mm}
  \begin{tikzpicture}
    \node[fill=accentColor!15, rounded corners=10pt, inner sep=5pt] {
      \color{accentColor}\bfseries\footnotesize CESS-UFF / SGIMP
    };
  \end{tikzpicture}

  \vspace{6mm}

  {\Large \textbf{\color{primaryColor} Sistema de Sinalização Embarcado}}\par
  \vspace{3mm}
  {\small \color{mutedText} ESP32, MicroPython, LEDs, Botão e OLED SSD1306}\par

  \vspace{8mm}

  \begin{tikzpicture}
    \node[fill=cardBg, draw=accentColor!30, rounded corners=8pt, inner sep=6pt, text width=\cardtextwidth, align=center] {
      {\small \textbf{Candidato:}}\\
      {\small \CandidateName}\\[3pt]
      {\small \textbf{Plataforma:} Wokwi (Simulador)}\\[3pt]
      {\small \textbf{Linguagem:} MicroPython}\\[3pt]
      {\small \textbf{Arquitetura:} Assíncrona (\texttt{asyncio})}
    };
  \end{tikzpicture}

  \vfill
\end{center}
\newpage

% ------------------------------------------------------------------------------
% SLIDE 2: Visão Geral do Projeto e Hardware
% (slides "Visão Geral" e "Hardware & Pinos" da versão anterior, fundidos)
% ------------------------------------------------------------------------------
\slideheader{Visão Geral \& Hardware}

\card{Objetivo do Projeto}{
  Desenvolver e simular um sistema embarcado responsivo com \textbf{ESP32} integrando tópicos de:
  \begin{itemize}
    \item Instrumentação
    \item Eletrônica
    \item Lógica de Programação
  \end{itemize}
}

\card{Placa-Alvo}{
  \textbf{Espressif ESP32-DevKitC V4}\\[2pt]
  Módulo recomendado: \textbf{ESP32-WROOM-32E}.\\[2pt]
  {\tiny (Evitar WROVER: usa internamente os GPIOs 16/17 para a PSRAM.)}
}

\alertcard{accentColor}{Conexões dos GPIOs}{
  \begin{tabular}{@{}ll@{}}
    \color{redAccent}\(\bullet\)\ \textbf{LED Vermelho:} & GPIO 2 (resistor 220~[$\Omega$]) \\
    \color{greenAccent}\(\bullet\)\ \textbf{LED Verde:} & GPIO 4 (resistor 220~[$\Omega$]) \\
    \(\bullet\)\ \textbf{Botão pulsador:} & GPIO 17 (\texttt{PULL\_DOWN}) \\
    \(\bullet\)\ \textbf{OLED SCL:} & GPIO 25 (I\textsuperscript{2}C \textit{clock}) \\
    \(\bullet\)\ \textbf{OLED SDA:} & GPIO 16 (I\textsuperscript{2}C \textit{data}) \\
  \end{tabular}
}
\newpage

% ------------------------------------------------------------------------------
% SLIDE 4: Funcionamento do Sistema
% ------------------------------------------------------------------------------
\slideheader{Funcionamento}

\alertcard{redAccent}{LED Vermelho (GPIO 2)}{
  Pisca continuamente a cada \textbf{500~[ms]}, em uma tarefa assíncrona própria.
}

\alertcard{greenAccent}{Resposta ao Botão (GPIO 17)}{
  \begin{tabular}{@{}lll@{}}
    \scriptsize\textbf{Botão} & \scriptsize\textbf{LED Verde} & \scriptsize\textbf{OLED} \\
    \scriptsize Solto (LOW) & \scriptsize Apagado & \scriptsize\texttt{"Boa sorte!"} \\
    \scriptsize Pressionado & \scriptsize Aceso & \scriptsize\texttt{"Consegui"} \\
  \end{tabular}
}

\card{Duas Tarefas, Nunca Bloqueantes}{
  Exclusivamente \textbf{\texttt{asyncio}} do MicroPython, sem superlaço (\textit{superloop}) nem \texttt{time.sleep()}: as duas tarefas acima rodam como \textbf{corrotinas} cooperativas, e nenhuma bloqueia a outra.
}

\alertcard{accentColor}{OLED: Atualização Dinâmica}{
  Atualizado apenas quando o estado \textbf{estável} do botão muda -- evita tráfego I\textsuperscript{2}C desnecessário e cintilação (\textit{flickering}).
}
\newpage

% ------------------------------------------------------------------------------
% SLIDE 5: Validação e Conclusão
% (slides "Validação" e "Conclusão" da versão anterior, fundidos)
% ------------------------------------------------------------------------------
\slideheader{Validação \& Conclusão}

\card{Estratégia de Testes Isolados}{
  Testes progressivos em \texttt{tests/}, do mais simples ao mais completo: LED vermelho (bloqueante e com \texttt{asyncio}), botão + LED verde, OLED I\textsuperscript{2}C (varredura e diagnóstico gráfico) e integração final em \texttt{main.py}.
}

\alertcard{greenAccent}{Requisitos Atendidos}{
  \begin{itemize}
    \item[\okmark] LED vermelho: ciclo de 500~[ms], contínuo e não bloqueante.
    \item[\okmark] LED verde: reflete o estado do botão.
    \item[\okmark] OLED: mensagem correta por estado, com atualização dinâmica.
    \item[\okmark] Execução concorrente via \texttt{asyncio}.
  \end{itemize}
}

\alertcard{primaryColor}{Links para Acesso}{
  Wokwi: \url{\WokwiURL}\\
  GitHub: \url{\GitHubURL}
}

\end{document}
